\begin{abstract}
Let \(\Phi\) be the classical Riemann, or de Bruijn--Newman, kernel. We prove
that every translation minor of \(\Phi\) of order at most four is strictly
positive at strictly ordered nodes, and a direct origin-confluent calculation
gives a negative order-five minor. Thus the exact global Polya-frequency order
is four. The computer-assisted input is confined to directed one-variable
inequalities: global positivity of the log-curvature quantities \(q,F_2\) and
of the fully confluent invariant \(\Cfour\), together with a five-derivative
origin certificate using exact rational arithmetic. The remainder is exact. A
quotient-Wronskian argument reduces arbitrary fourth-order minors to
\(\partial_\xi\Psi<0\). In the increasing curvature coordinate, its numerator
is
\[
 \delta\Lambda\int_p^w
 W(t)\frac{\Cfour(t)}{Q(t)^6\kappa(t)^2}\dd t.
\]
Here \(W\) is the difference of two explicitly normalized triangular CDFs.
Its strict positivity follows from a one-crossing density ratio. All
load-bearing algebraic reductions are displayed in the text and appendices;
the repository scripts provide replay rather than substitute derivations.
We then construct an explicit positive even Schwartz kernel which is itself
strictly \(\PF_4\) but whose entire Fourier transform has nonreal zeros. Thus
continuous \(\PF_3\) and \(\PF_4\) membership alone do not force Fourier
real-rootedness. The results are finite-order and have no implication for the
Riemann Hypothesis.
\end{abstract}
